% TODO make hw stylesheet
\documentclass[11pt]{scrartcl}
\usepackage[a4paper, total={6in, 8in}]{geometry}
\usepackage[usenames,dvipsnames,svgnames]{xcolor}
\usepackage[shortlabels]{enumitem}
\usepackage[framemethod=TikZ]{mdframed}
\usepackage{amsmath,amssymb,amsthm}
\usepackage[colorlinks]{hyperref}
\usepackage{multicol}
\usepackage{tabularx}

\hypersetup{
	colorlinks=true,
	linkcolor=blue,
	filecolor=magenta,      
	urlcolor=magenta,
	pdftitle={MAT 320 HW}
}

\usepackage{microtype}
\usepackage{mathtools}
\usepackage[headsepline]{scrlayer-scrpage}
\usepackage{thmtools}
\usepackage[textsize=footnotesize]{todonotes}

\mdfdefinestyle{mdbluebox}{roundcorner=10pt,innerbottommargin=9pt,
	linecolor=blue,backgroundcolor=TealBlue!5,}
\declaretheoremstyle[headfont=\sffamily\bfseries\color{MidnightBlue},
mdframed={style=mdbluebox},]{thmbluebox}

\mdfdefinestyle{mdbloobox}{frametitlefont=\bfseries,innerbottommargin=8pt,
	nobreak=true,backgroundcolor=TealBlue!5,linecolor=blue,}
\declaretheoremstyle[headfont=\bfseries\color{MidnightBlue},
mdframed={style=mdbloobox},headpunct={\\[3pt]},postheadspace=0pt,]{thmbloobox}

\mdfdefinestyle{mdredbox}{frametitlefont=\bfseries,innerbottommargin=8pt,
	nobreak=true,backgroundcolor=Salmon!5,linecolor=RawSienna,}
\declaretheoremstyle[headfont=\bfseries\color{RawSienna},
mdframed={style=mdredbox},headpunct={\\[3pt]},postheadspace=0pt,]{thmredbox}

\mdfdefinestyle{mdgreenbox}{linecolor=ForestGreen,backgroundcolor=ForestGreen!5,
	linewidth=2pt,rightline=false,leftline=true,topline=false,bottomline=false,}
\declaretheoremstyle[headfont=\bfseries\sffamily\color{ForestGreen!70!black},
mdframed={style=mdgreenbox},headpunct={ --- },]{thmgreenbox}

\mdfdefinestyle{mdorangebox}{linecolor=RedOrange,backgroundcolor=RedOrange!5,
	linewidth=2pt,rightline=false,leftline=true,topline=false,bottomline=false,}
\declaretheoremstyle[headfont=\bfseries\sffamily\color{RedOrange!70!black},
mdframed={style=mdorangebox},headpunct={ --- },]{thmorangebox}

\mdfdefinestyle{mdblackbox}{linecolor=black,backgroundcolor=RedViolet!5!gray!5,
	linewidth=3pt,rightline=false,leftline=true,topline=false,bottomline=false,}
\declaretheoremstyle[mdframed={style=mdblackbox}]{thmblackbox}

\declaretheoremstyle[spaceabove=6pt,spacebelow=6pt]{basehead}

\declaretheorem[style=basehead,name=Problem]{problem}
\declaretheorem[style=thmredbox,name=Problem,numbered=no]{problem*}
\declaretheorem[style=thmbloobox,name=Setup]{setup} % for config geo
\declaretheorem[style=thmredbox,name=Example,sibling=problem]{example}
\declaretheorem[style=thmredbox,name=$\bigstar$ Problem,sibling=problem]{niceprob}
\declaretheorem[style=thmbluebox,name=Theorem,sibling=problem]{theorem}
\declaretheorem[style=thmbluebox,name=Corollary,sibling=theorem]{corollary}
\declaretheorem[style=thmbluebox,name=Proposition,sibling=theorem]{prop}
\declaretheorem[style=thmbluebox,name=Lemma,numberwithin=problem]{lemma}
\declaretheorem[style=thmbluebox,name=Theorem,numbered=no]{theorem*}
\declaretheorem[style=thmbluebox,name=Lemma,numbered=no]{lemma*}
\declaretheorem[style=thmgreenbox,name=Claim,numberwithin=problem]{claim}
\declaretheorem[style=thmgreenbox,name=Claim,numbered=no]{claim*}
\declaretheorem[style=thmblackbox,name=Remark,numberwithin=problem]{remark}
\declaretheorem[style=thmblackbox,name=Remark,numbered=no]{remark*}
\declaretheorem[style=thmgreenbox,name=Definition,sibling=theorem]{definition}
\declaretheorem[style=thmgreenbox,name=Definition,numbered=no]{definition*}
\declaretheorem[style=thmorangebox,name=Warning,sibling=theorem]{warning}
\declaretheorem[style=thmorangebox,name=Warning,numbered=no]{warning*}
\declaretheorem[style=thmorangebox,name=Note,numbered=no]{note}
\declaretheorem[style=thmblackbox,name=Exercise,sibling=problem]{exercise}
\declaretheorem[style=thmblackbox,name=Exercise,numbered=no]{exercise*}
\declaretheorem[style=thmblackbox,name=Question,sibling=problem]{question}
\declaretheorem[style=thmblackbox,name=Question,numbered=no]{question*}

\addtolength{\textheight}{3.14cm}
\providecommand{\ol}{\overline}
\providecommand{\ul}{\underline}
\providecommand{\eps}{\varepsilon}
\providecommand{\half}{\frac{1}{2}}
\providecommand{\dang}{\measuredangle} %% Directed angle
\providecommand{\CC}{\mathbb C}
\providecommand{\FF}{\mathbb F}
\providecommand{\NN}{\mathbb N}
\providecommand{\QQ}{\mathbb Q}
\providecommand{\RR}{\mathbb R}
\providecommand{\ZZ}{\mathbb Z}
\providecommand{\dg}{^\circ}
\providecommand{\ii}{\item}
\providecommand{\setdiff}{\smallsetminus}
\providecommand{\alert}[1]{{\sffamily\textbf{\textcolor{blue}{#1}}}}
\providecommand{\inv}{^{-1}}
\providecommand{\mb}[1]{\mathbf{#1}}

\reversemarginpar

\newenvironment{soln}{\begin{proof}[Solution]}{\end{proof}}
\newenvironment{walkthrough}{\noindent \textbf{Walkthrough.}}{\medskip}
\newlist{walk}{enumerate}{3}
\setlist[walk]{label=\bfseries (\alph*)}

\providecommand{\pow}{\operatorname{Pow}}
\providecommand{\vocab}[1]{{\textbf{\textcolor{ForestGreen}{#1}}}}
\providecommand{\lcm}{\operatorname{lcm}}
\providecommand{\abs}[1]{\left|#1\right|}

\usepackage{asymptote}
\begin{asydef}
	defaultpen(fontsize(10pt));
	unitsize(1cm);
	size(8cm); // set a reasonable default
	usepackage("amsmath");
	usepackage("amssymb");
	settings.tex="pdflatex";
	settings.outformat="pdf";
	// Replacement for olympiad+cse5 which is not standard
	import geometry;
	// recalibrate fill and filldraw for conics
	void filldraw(picture pic = currentpicture, conic g, pen fillpen=defaultpen, pen drawpen=defaultpen)
	{ filldraw(pic, (path) g, fillpen, drawpen); }
	void fill(picture pic = currentpicture, conic g, pen p=defaultpen)
	{ filldraw(pic, (path) g, p); }
	// some geometry
	pair foot(pair P, pair A, pair B) { return foot(triangle(A,B,P).VC); }
	pair orthocenter(pair A, pair B, pair C) { return orthocentercenter(A,B,C); }
	pair centroid(pair A, pair B, pair C) { return (A+B+C)/3; }
	// cse5 abbreviations
	path CP(pair P, pair A) { return circle(P, abs(A-P)); }
	path CR(pair P, real r) { return circle(P, r); }
	pair IP(path p, path q) { return intersectionpoints(p,q)[0]; }
	pair OP(path p, path q) { return intersectionpoints(p,q)[1]; }
	path Line(pair A, pair B, real a=0.6, real b=a) { return (a*(A-B)+A)--(b*(B-A)+B); }
	// cse5 more useful functions
	picture CC() {
		picture p=rotate(0)*currentpicture;
		currentpicture.erase();
		return p;
	}
	pair MP(Label s, pair A, pair B = plain.S, pen p = defaultpen) {
		Label L = s;
		L.s = "$"+s.s+"$";
		label(L, A, B, p);
		return A;
	}
	pair Drawing(Label s = "", pair A, pair B = plain.S, pen p = defaultpen) {
		dot(MP(s, A, B, p), p);
		return A;
	}
	path Drawing(path g, pen p = defaultpen, arrowbar ar = None) {
		draw(g, p, ar);
		return g;
	}
\end{asydef}

\usepackage{mathrsfs}

\ihead{\footnotesize MAT 320 HW08}
\ohead{\footnotesize November 1, 2024}

\title{\vspace{-2em}MAT 320 HW08}
\author{\large Aditya Pahuja}
\date{\large Due 11/1}
\begin{document}
\maketitle

\begin{problem}
	Let $(X, d)$ be a metric space and let $E$ be a subset of $X$.
	Show that if $U\subseteq X$ is open then
	$U\cap E = \varnothing$ if and only if $U\cap\ol E = \varnothing$.
\end{problem}

\begin{soln}
	If $U\cap\ol E = \varnothing$ then $U\cap E = \varnothing$ is trivial
	because $E\subseteq\ol E$, as we showed in last week's homework.
	
	Conversely, suppose for the sake of contradiction that
	$U\cap E = \varnothing$ but $U\cap \ol E$ is nonempty.
	Let $x$ be in $U\cap \ol E$. Since $x$ is a limit point of $E$,
	we can find a sequence $(x_n)$ of points in $E$
	whose limit is $x$. In particular, for every $\eps > 0$,
	$x_n\in B(x,\eps)$ eventually.
	Because $U$ is open, we can pick $\eps$ such that $B(x, \eps)\subseteq U$.
	However, there exists an $n$ such that $x_n\in B(x,\eps)$,
	which means $x_n$ is in both $U$ and $E$; this is contradicts
	$U\cap E = \varnothing$, so $U\cap\ol E$ is empty.
\end{soln}

\begin{problem}
	Let $(X, d)$ be a metric space.
	Show that $(X, d)$ is compact if and only if
	for any collection of closed sets $\{F_i\subseteq X : i\in I\}$
	such that $\bigcap_{i\in I} F_i = \varnothing$,
	there is a finite subcollection $\{F_i\subseteq X:i\in J\}$
	such that $\bigcap_{i\in J} F_i = \varnothing$.
\end{problem}
\begin{soln}
	Let $\mathcal U = \{U_i : i\in I\}$ be any open cover of $X$.
	Letting $F_i = X\setdiff U_i$, we see that
	\[\bigcap_{i\in I} F_i = \bigcap_{i\in I} X\setdiff U_i
	= X\setdiff \bigcup_{i\in I} U_i = \varnothing.\]
	This means there exists a finite set of indices $i_1$, \dots, $i_n$
	such that
	\[\varnothing = \bigcap_{k=1}^n F_{i_k} =
	\bigcap_{k=1}^n X\setdiff U_{i_k} = X\setdiff\bigcup_{k=1}^n U_{i_k}
	\iff X = \bigcup_{k=1}T^n U_{i_k}.\]
	That is, given any open cover $\mathcal U$ of $X$,
	we have extracted a finite collection of $\mathcal U$
	(namely $\{U_{i_k} : 1\le k\le n\}$)
	that is an open subcover, which is exactly what it means
	for $X$ to be compact.
\end{soln}

\begin{problem}
	Prove directly from the definition of compact set that the set
	\[K = \{0\}\cup\{1/n\in\RR: n\in\NN\}\]
	is a compact subset of $\RR$.
\end{problem}

\begin{soln}
	Let $\mathcal U = \{U_i : i\in I\}$ be an open cover of $X$.
	One of these $U_i$ must contain $0$; call this set $U$.
	By definition of openness, there exists an $\eps > 0$
	such that $B(0, \eps)\subseteq U$. Let $N$ an integer
	such that $\eps > \frac 1N$.
	If $n < N$, then $\abs{\frac 1n - 0} < \eps$, meaning
	$\frac 1n \in B(0,\eps) \subseteq U$.
	If $n \ge N$, then there exist indices $i_1$, \dots, $i_N$
	such that $U_{i_k}$ contains $\frac 1k$
	because each element of $K$ is contained in
	at least one open set in $\mathcal U$.
	These two cases encompass all positive integers, so
	we have found $N+1$ (i.e. finitely many) sets in $\mathcal U$ such that
	\[U\cup U_{i_1} \cup U_{i_2} \cup\cdots\cup U_{i_N} \supset K.\]
	That is, we have extracted a finite subcover of $\mathcal U$,
	hence $K$ is compact.
\end{soln}

\begin{problem}
	Show that $(0,1)$ is not compact by showing an example of
	an open cover admitting no open subcover.
	Do you know other ways of showing that $(0,1)$ is not compact?
	Challenge: find as many different proofs as possible.
\end{problem}

\begin{proof}[Solution 1]
	Consider the open cover
	\[\mathcal U = \{U_\alpha = (0, \alpha) : \alpha\in(0,1)\}.\]
	Then, given any finite collection of indices
	$\alpha_1,\,\alpha_2,\,\cdots,\,\alpha_n\in (0,1)$
	that are WLOG ordered such that $\alpha_1 \le \cdots \le \alpha_n$, we get
	\[U_{\alpha_1}\cup\cdots\cup U_{\alpha_n} = (0, \alpha_n)\]
	because $U_\alpha \subseteq U_\beta$ if $\alpha\le\beta$.
	However, this union is clearly not equal to $(0,1)$,
	so this finite subcollection of $\mathcal U$ is not a cover of $(0,1)$.
\end{proof}

\begin{proof}[Solution 2]
	We showed in class that compact sets are always closed.
	However, $(0,1)$ is not closed, so it is not compact.
\end{proof}

\begin{proof}[Solution 3]
	Alternatively, we showed in class that a subset of $\RR^n$ is compact
	if and only if it is closed and bounded.
	Again, $(0,1)$ is not closed, so it cannot be compact.
\end{proof}

\begin{proof}[Solution 4]
	Since compactness implies sequential compactness,
	we can prove that $(0,1)$ is not compact
	by proving that it is not sequentially compact.
	Indeed, the sequence $x_n = \frac 1n$ for $n > 1$
	converges to $0$ in $\RR$, so any subsequence
	also converges to $0$ in $\RR$.
	Therefore, no subsequence of $(x_n)$ converges in $(0,1)$,
	so $(0,1)$ is not sequentially compact, hence not compact.
\end{proof}

\begin{problem}
	Show that if $C\subseteq\RR^n$ is bounded then $K = \ol C$ is compact.
\end{problem}

\begin{soln}
	We note that $\ol C$ is bounded:
	\begin{claim}
		$\operatorname{diam}C = \operatorname{diam}\ol C$.
	\end{claim}
	\begin{proof}
		Since $C\subseteq\ol C$,
		\[\operatorname{diam} C \le \operatorname{diam}\ol C.\]
		
		To show the opposite inequality, we prove equivalently that
		\[\operatorname{diam}\ol C \le \operatorname{diam} C + \eps\]
		for every $\eps > 0$.
		Indeed, let $x,\,y\in\ol C$. Then, $x$ and $y$ are limit points of $C$,
		so we can find $x',\,y'\in C$ such that
		$d(x,x')$ and $d(y, y')$ are both less than $\frac\eps 2$.
		The triangle inequality then says that
		\[d(x, y) \le d(x, x') + d(x', y') + d(y', y)
		< \frac\eps 2 + \operatorname{diam}C + \frac\eps2
		= \operatorname{diam} C + \eps.\]
		Thus $\operatorname{diam} C + \eps$ is an upper bound
		for $d(x, y)$, meaning it is at least as large as
		$\operatorname{diam} C$, as desired.
	\end{proof}
	
	Since $\operatorname{diam}C$ is finite,
	so is $\operatorname{diam}\ol C$, hence $\ol C$ is bounded.
	Moreover, $\ol C$ is closed, as shown in the previous homework.
	Thus, $\ol C$ a closed and bounded subset of $\RR^n$,
	so, as we proved in class, it is compact.
\end{soln}


\begin{problem}
	A function $f\colon X\to Y$ between metric spaces is called
	a \vocab{Lipschitz} function if there exists a constant $L > 0$ such that
	\[d_Y(f(x), f(y)) < L\cdot d_X(x, y)\]
	for all $x,\,y\in X$.
	Show that Lipschitz functions are uniformly continuous.
\end{problem}

\begin{soln}
	Given $\eps > 0$, let $\delta = \frac{\eps}{L}$. Then,
	\[d_X(x,y) < \delta\implies d_Y(f(x),f(y)) <
	L\cdot d_X(x,y) < L\cdot\delta = \eps,\]
	which is what it means for $f$ to be uniformly continuous.
\end{soln}

\begin{problem}
	Consider the set $\QQ\subseteq\RR$ with the induced metric
	\[d(p, q) = |p - q|.\]
	Show that $C = \{q\in\QQ : |q| < \sqrt2\}$
	is closed (in $\QQ$) and bounded, but not compact.
\end{problem}

\begin{soln}
	We can write $\QQ\setdiff C$ as
	\[\QQ\cap((-\infty, -\sqrt2)\cup(\sqrt2,\infty)).\]
	This is an intersection of $\QQ$ with a set that is open in $\RR$
	(because it is a union of open intervals),
	so $\QQ\setdiff C$ must be open in $\QQ$.
	Then, the set's complement in $\QQ$, namely $C$ is closed in $\QQ$.
	
	Moreover, $C$ is bounded because, for any $x,y\in C$
	with WLOG $x \ge y$, we have $-\sqrt 2 < y \le x < \sqrt2$, so
	\[|x - y| = x - y < \sqrt2 - (-\sqrt 2) = 2\sqrt2,\]
	hence the diameter of $C$ is bounded above by $2\sqrt2$.
	
	To see that $C$ is not compact, it suffices to prove that
	it is not sequentially compact (taking the contrapositive of
	``compact implies sequentially compact'').
	Indeed, we can take any sequence of rationals
	less than $\sqrt2$ that converges to $\sqrt 2$ (in $\RR$);
	clearly, this sequence has no subsequential limits in $\QQ$.
	One example of such a sequence is
	\[a_n = \frac{1}{10^n}\cdot\lfloor 10^n\sqrt2\rfloor,\]
	in which $\lfloor x\rfloor$ is the integer part of $x$,
	so $a_n$ is the truncation of the decimal representation of $2$
	after $n$ digits past the decimal point
	(so $a_0 = 1$, $a_1 = 1.4$, $a_2 = 1.41$, and so on).
	This sequence is always in $C$ and converges to $\sqrt 2$
	because
	\[0 \le \sqrt2 - a_n = \frac{10^n\sqrt2 - \lfloor10^n\sqrt2\rfloor}{10^n}
	< \frac{1}{10^n} \to 0,\]
	as the numerator, the fractional part of $10^n\sqrt2$,
	is always nonnegative but less than $1$.
\end{soln}

\begin{problem}
	Let $(X, d)$ be a compact metric space and
	let $\mathcal U = \{U_i\subseteq X : i\in I\}$ be an open cover.
	Show that there exists $\delta > 0$ such that
	any subset of $X$ having diameter less than $\delta$
	is contained in at least one of the open sets in $\mathcal U$.
\end{problem}

\begin{soln}
	Since $X$ is compact, it is covered by finitely many of the $U_i$;
	call this finite subcover $U_1$, $U_2$, \dots, $U_n$.
	For each $k = 1,\,\dots,\, n$, define the function
	$f_k\colon X\to\RR$ as follows:
	\begin{itemize}
		\ii If $x\in U_k$, then
		\[f_k(x) = \sup\{r : B(x, r)\subseteq U_k\}.\]
		Note that $U_k$ being open means that this value is positive.
		\ii If $x\notin U_k$, then $f_k(x) = 0$.
	\end{itemize}
	Thus, $f_k(x) = r$ means that
	any open ball centered at $x$ of radius less than $r$ is contained in $U_k$
	(with $r = 0$ meaning no open ball centered at $x$ is a subset of $U_k$).

	\begin{claim}
		Each $f_k$ is a continuous function.
	\end{claim}

	\begin{proof}
		Let $\eps > 0$ and $x\in X$. Suppose $d(x, y) < \eps$.
		Then, we have four cases:
		\begin{enumerate}[(a)]
			\ii Both $x$ and $y$ are in $U_k$. Then, letting $r = f_k(x)$,
			\[B(y,r-\eps) \subseteq B(x,r)\subseteq B(y,r+\eps)\]
			shows that $|f_k(y) - r| < \eps$.
			\ii $x\in U_k$ and $y$ is not.
			In this case, $y\in B(x, \eps)$ means that
			$B(x, \eps)\not\subseteq U_k$, so $f_k(x) \le d(x, y) < \eps$.
			Thus $|f_k(x) - f_k(y)| = |f_k(x)| < \eps$.
			\ii $y\in U_k$ and $x$ is not.
			This is entirely symmetric to the case above:
			we deduce that $f_k(y) < \eps$, so
			$|f_k(x) - f_k(y)| = |f_k(y)| < \eps$.
			\ii Neither $x$ nor $y$ are in $U_k$.
			In this case we have $|f_k(x) - f_k(y)| = 0 < \eps$.
		\end{enumerate}
		This shows that, if $d(x, y) < \eps$, then
		\[d_{\RR}(f_k(x), f_k(y)) = |f_k(x) - f_k(y)| < \eps,\]
		so $f_k$ is continuous.
	\end{proof}
	
	\begin{claim}
		The function
		\[F(x) = \max\{f_1(x), \dots, f_n(x)\}\]
		is continuous.
	\end{claim}

	\begin{proof}
		We prove this by induction.
		The idea is that, for $n = 2$,
		\[\max\{f_1(x), f_2(x)\} =
		\frac{|f_1(x) + f_2(x)| + |f_1(x) - f_2(x)|}{2},\]
		which is a sum of a composition of continuous functions
		and is thus continuous.
		Then, if the statement holds for some $n$,
		it holds for $n + 1$ because
		\[\max\{f_1(x), \dots, f_n(x), f_{n+1}(x)\} =
		\max\{\max\{f_1(x), \dots, f_n(x)\}, f_{n+1}(x)\},\]
		and we can use the two-function case again.
	\end{proof}

	Observe that $F(x) = r$ means that
	any open ball centered at $x$ with radius less than $r$
	will be completely contained within one of the $U_k$.
	Moreover, $F(x) > 0$ for all $x$ because each $x$
	is in at least one of the $U_k$, so $f_k(x) > 0$
	for at least one $k$ regardless of the value of $x$.
	Finally, $F$ is a continuous function
	whose domain is compact, which means that
	$F(X)$ is a closed subset of $\RR$.
	Crucially, this means that $F$ attains a minimum value $\delta > 0$.
	
	To finish, let $E$ have diameter less than $\delta$.
	Then, $E\subseteq B(x,\delta)$ for any $x\in E$,
	so there exists a $k\in\{1,\dots,n\}$ such that
	$$E\subseteq B(x,\delta)\subseteq U_k.\qedhere$$
\end{soln}

\begin{problem}
	Prove that a function $f\colon[0,1]\to\RR$ is continuous
	if and only if its graph
	\[\Gamma_f = \{(x, f(x)) : x\in[0,1]\}.\]
	is a compact subset of $\RR^2$.
	Give an example of a function $f\colon[0,1]\to\RR$
	that is discontinuous with graph $\Gamma_f$ which is
	closed but not compact.
\end{problem}

\begin{soln}
	Define $\pi\colon[0,1]\to\Gamma_f$ by
	\[\pi(x) = (x, f(x)).\]
	\begin{claim}
		$\pi$ is continuous if and only if $f$ is continuous.
	\end{claim}
	\begin{proof}
		First, suppose $f$ is continuous, and let $c\in[0,1]$ and $\eps > 0$.
		By continuity, there exists a $\delta_1$ such that
		\[|x - c| < \delta_1 \implies |f(x) - f(c)| < \eps_1\]
		where $0 < \eps_1 < \eps$.
		Of course,
		\[|x - c| + |f(x) - f(c)| < \eps_1 + \delta_1,\]
		i.e. $\pi(x)$ and $\pi(c)$ have a distance less than $\eps_1+\delta_1$
		in the taxicab metric. Now, we can pick
		\[\delta < \min\{\eps - \eps_1, \delta_1\}\]
		so that
		\[|x - c| < \delta \implies |x - c| + |f(x) - f(c)| <
		\eps_1 + \delta < \eps_1 + \eps - \eps_1 = \eps.\]
		This is exactly what it means for $\pi$ to be continuous,
		again recalling that the term that we have bounded above by $\eps$
		is just the distance between $\pi(x)$ and $\pi(c)$.
		
		For the other direction, if $\pi$ is continuous, then
		for each $\eps > 0$ there is a $\delta$ such that
		\[|x - c| < 0 \implies |f(x) - f(c)|
		< d_{\text{taxicab}}(\pi(x),\pi(c)) < \eps,\]
		so we get continuity of $f$ for free.
	\end{proof}

	Now, if $f$ is continuous, then $\pi$ is continuous, meaning
	$\pi([0,1]) = \Gamma_f$ is compact because it is
	the continuous image of a compact set.
	On the other hand, if $\Gamma_f$ is compact,
	consider $\pi\inv$, which maps $(x, f(x))\mapsto x$.
	This is uniformly continuous because for each $\eps > 0$,
	we can take $\delta = \eps$ to get
	\[d_{\text{taxicab}}((x, f(x)), (y, f(y))) =
	|x - y| + |f(x) - f(y)| < \eps = \delta \implies
	|x - y| < \eps\]
	for any two points $(x, f(x))$ and $(y, f(y))$ in $\Gamma_f$.
	By part (e) of Problem 14, $\pi\inv$ has a continuous inverse
	(it is bijective because $\pi\inv$ has a inverse),
	so $\pi$ is continuous and thus $f$ is continuous.
	
	An example of $f$ such that $\Gamma_f$ is closed but not compact:
	take $f(x) = \frac{1}{x}$ for $x\ne 0$ and $f(0) = 0$.
	The $\pi\inv$ map from earlier is still continuous,
	so the pre-image $\pi([0,1]) = \Gamma_f$ is closed
	because $[0,1]$ is closed. However, clearly the graph is unbounded
	because $\frac 1x$ is unbounded on $[0,1]$, so
	$\Gamma_f$ cannot be compact.
\end{soln}

\setcounter{problem}{10}

\begin{problem}
	Show that in $\ell^2$ the set
	\[B = \left\{\mb x = (x_1, x_2, x_3, \dots) \in\ell^2 :
	\sum_{i=1}^{\infty}x_i^2 \le 1\right\}\]
	is closed and bounded but not compact.
\end{problem}

\begin{soln}
	To show that $B$ is closed, we use a little trick.
	Denoting the norm of $\mb x$ as $\|x\|$, we can see that
	$f\colon \ell^2\to \RR$ given by $f(x) = \|x\|$
	is continuous because, for each $\eps > 0$,
	we can take $\delta = \eps$ to get
	\[d(\mb x, \mb y) < \eps \implies
	\abs{\|\mb x\| - \|\mb y\|} \le \|\mb x - \mb y\| = d(\mb x,\mb y) < \eps,\]
	where the inequality of norms follows from the triangle inequality:
	if WLOG $\|\mb x\| \ge \|\mb y\|$, then
	\[\|\mb x\| = d(\mb x,\mb 0) \le d(\mb x, \mb y) + d(\mb y, \mb 0)
	= \|\mb x - \mb y\| + \|\mb y\| \implies
	\|\mb x\| - \|\mb y\| \le \|\mb x-\mb y\|.\]
	Since $B$ is by definition the set of all $\mb x$ such that $\mb x \le 1$,
	$f(B)\subseteq[0,1]$, and, for each $\lambda\in[0,1]$,
	$\mb x = (\lambda, 0,0,\dots)\in B$ means that
	\[f(\mb x) = \sqrt{\lambda^2} = |\lambda| = \lambda \in f(B).\]
	That is, $[0,1]\subseteq f(B)$ too, which proves that
	$[0,1] = f(B)$ is a closed set in $\RR$.
	Since $f$ is continuous, the pre-image of a closed set
	is also closed, so $f\inv([0,1]) = B$ is closed in $\ell^2$.
	
	Additionally, $B$ is bounded because for any $\mb x$ and $\mb y$ in $B$,
	we can use the fact that $\|\mb x\| \le 1$ and $\|\mb y\|\le 1$ to say that
	\[d(\mb x, \mb y) = \|\mb x - \mb y\|\le \|\mb x\| + \|-\mb y\|
	= \|\mb x\| + \|\mb y\| \le 2,\]
	so any two points in $B$ have a distance at most two,
	hence $\operatorname{diam} B \le 2 < +\infty$.
	
	\begin{remark}
		We are taking for granted here the vector space structure of $\ell^2$
		that is established in Problem 12.
	\end{remark}

	Finally, we show that $B$ is not compact by proving that
	it is not sequentially compact.
	Indeed, the sequence
	\[\mb x_n =
	(\underbrace{0, 0, \dots, 0}_{n-1\text{ zeroes}}, 1, 0, 0, \dots)\]
	where the $n$th coordinate of $\mb x_n$ is $1$
	and the rest are zero does not have any convergent subsequences
	because it doesn't even have any Cauchy subsequences, since
	\[d(\mb x_m, \mb x_n) = \sqrt2\]
	for all $m$ and $n$ with $m\ne n$.
\end{soln}

\begin{problem}
	Two infinite vectors $\mb x = (x_1, x_2, \dots)$
	and $\mb y = (y_1, y_2, \dots)$ can be summed:
	\[\mb x + \mb y = \{x_1 + y_1, x_2 + y_2, \dots\}\]
	and multiplied by a constant $\lambda\in\RR$:
	\[\lambda\mb x = (\lambda x_1, \lambda x_2, \dots).\]
	Complete the following tasks:
	\begin{enumerate}[(a)]
		\ii Show that if $\mb x$ and $\mb y$ are in $\ell^2$
		then so are $\mb x + \mb y$ and $\lambda\mb x$
		(i.e. $\ell^2$ is a vector space
		with these addition and scalar multiplication operations).
		\ii Give a definition for what it means for a series
		\[\sum_{i=1}^{\infty} \mb x_i\]
		of vectors $\mb x_1,\,\mb x_2,\,\dots\in\ell^2$ to converge.
		\ii Use your definition to give meaning to infinite linear combinations
		\[\mb x = \sum_{i=1}^\infty \lambda_i\mb x_i.\]
		\ii Show that any vector $\mb x\in\ell^2$ can be written as
		an infinite linear combination of the ``standard basis vectors''
		$\mb e_1 = (1, 0, 0, 0, \dots)$, $\mb e_2 = (0, 1, 0, 0, \dots)$,
		$\mb e_3 = (0, 0, 1, 0, \dots)$, and so on.
	\end{enumerate}
\end{problem}

\begin{soln}
	For (a), if $\mb x$ and $\mb y$ are in $\ell^2$, then
	\[\sum_{i=1}^\infty (x_i + y_i)^2 =
	\sum_{i=1}^\infty x_i^2 + y_i^2 + 2x_iy_i
	= \|\mb x\|^2 + \|\mb y\|^2 + 2\sum_{i=1}^\infty x_iy_i,\]
	which converges because $|x_iy_i|$ can be bounded termwise by
	$\frac{x_i^2 + y_i^2}{2}$ (e.g. by reducing the bound to the equivalent
	$(|x_i| - |y_i|)^2 \ge 0$).
	In other words $\mb x + \mb y$ is in $\ell^2$
	because $\sum x_iy_i$ converges absolutely, hence it converges.
	Additionally,
	\[\sum_{i=1}^\infty (\lambda x_i)^2 = \lambda^2 \sum_{i=1}^\infty x_i^2\]
	converges when $\mb x\in\ell^2$, so $\lambda\mb x\in\ell^2$
	for any real $\lambda$ too.
	
	For (b) we can define convergence of series in the obvious way:
	we define
	\[\sum_{i=1}^\infty \mb x_i = \lim_{n\to\infty}\sum_{i=1}^n \mb x_i\]
	to be the limit of the sequence of the partial sums of $\mb x_i$,
	if it converges.
	
	For (c) we then define the infinite linear combination
	to just be the value of the respective infinite sum?
	
	Finally, for (d), I claim that we can write
	\[\mb x = \sum_{i=1}^\infty x_i\mb e_i.\]
	Indeed, the $n$th partial sum of $x_i\mb e_i$ is
	$\mb x_n = (x_1, x_2, \dots, x_n, 0, 0, \dots)$, so
	\[\|\mb x - \mb x_n\| = \left\|\sum_{i=n+1}^\infty x_i\mb e_i\right\|
	= \sum_{i=n+1}^\infty x_i^2 = \|\mb x\|^2 - \sum_{i=1}^n x_i^2.\]
	Taking the limit as $n$ goes to infinity,
	\[\|\mb x\|^2 - \sum_{i=1}^n x_i^2 \to 0,\]
	so $\|\mb x - \mb x_n\|$ becomes arbitrarily small as $n$ grows large,
	meaning $\mb x_n$ converges to $\mb x$.
\end{soln}

\begin{problem}
	A function $\varphi\colon\ell^2\to\RR$ is called \vocab{linear} if
	\[\varphi(\mb x + \mb y) = \varphi(\mb x) + \varphi(\mb y)\quad
	\text{and}\quad
	\varphi(\lambda\mb x) = \lambda\varphi(\mb x)\]
	for all $\mb x,\,\mb y\in\ell^2$.
	Show that if $\varphi\colon\ell^2\to\RR$ is linear
	then the following statements are equivalent:
	\begin{enumerate}[(i)]
		\ii $\varphi$ is continuous;
		\ii $\varphi$ is uniformly continuous;
		\ii $\varphi$ is Lipschitz; and
		\ii $\varphi$ has bounded norm, meaning
		\[\sup\left\{|\varphi(\mb x)| : \mb x\in\ell^2\text{ such that }
		\sum_{i=1}^\infty x_i^2 \le 1\right\} < +\infty.\]
	\end{enumerate}
\end{problem}

\begin{soln}
	We know that
	\[\mathrm{(iii)}\implies \mathrm{(ii)}\implies \mathrm{(i)}.\]
	We will prove additionally that (i) implies (iv) and (iv) implies (iii).
	
	First we show (i) implies (iv).
	Suppose for the sake of contradiction that
	$\varphi$ does not have a bounded norm.
	
	\begin{claim}
		$|\varphi(\mb x)|$ is unbounded on every open ball centered at $\mb 0$.
	\end{claim}
	\begin{proof}
		Obviously $\varphi$ being bounded on the closed unit ball
		means it is bounded on the open unit ball. Define
		\[L = \sup\{|\varphi(\mb x)| : \|\mb x\| < 1\}\]
		and let $\eps > 0$. Then, we have
		\[\sup\{|\varphi(\mb y)| : \|\mb y\| < \eps\}
		= \sup\{|\varphi(\mb x\cdot\eps)| : \|\mb x\| < 1\}
		= \sup\{\eps\cdot |\varphi(\mb x)| : \|\mb x\| < 1\}
		= L\eps,\]
		so $\varphi$ is bounded on the open ball of radius $\eps$
		for any $\eps > 0$, as desired.
	\end{proof}

	On the other hand, let $U$ be the pre-image of the open interval $(-1,1)$.
	Since $\varphi$ is continuous, $U$ is open in $\ell^2$. Note that
	\[\varphi(\mb 0) = \varphi(2\cdot\mb 0) = 2\cdot\varphi(\mb 0),\]
	so $\varphi(\mb 0) = 0$, hence $\mb 0\in U$. In particular,
	$U$ also contains some open ball $B$ containing $\mb 0$.
	That is, for all $\mb x\in B$, $\varphi(\mb x) < 1$,
	which means $\varphi$ is bounded on $B$, contradicting the claim.
	
	To show that (iv) implies (iii), we need to prove that
	if $\varphi$ has a bounded norm then there exists a constant $L$ such that
	\[|\varphi(\mb x - \mb y)| =
	|\varphi(\mb x) - \varphi(\mb y)| \le L\cdot\|\mb x - \mb y\|.\]
	This is obvious when $\mb x = \mb y$, so assume otherwise.
	Then, letting $\mb u = \mb x - \mb y$, we want to show that
	the function $f\colon\ell^2\setdiff \mb 0\to \RR$
	\[f(\mb u) = \frac{|\varphi(\mb u)|}{\|\mb u\|}\]
	is bounded as $\mb u$ varies across the domain of $f$.
	Noting that $f(\lambda\mb u) = f(\mb u)$ for any nonzero real $\lambda$,
	we can assume without loss of generality that $\|\mb u\| = 1$.
	The result is then immediate by our assumption that
	$\varphi$ has bounded norm, since
	$f(\mb u) = |\varphi(\mb u)|$ is bounded on
	$\{\mb u\in\ell^2 : \|\mb u\|=1\}.$
\end{soln}

\begin{problem}
	Two metric spaces $X$ and $Y$ are \vocab{topologically equivalent}
	or \vocab{homeomorphic} if there exists a \vocab{homeomorphism}
	between them, which is a bijection $f\colon X\to Y$ such that
	both $f$ and $f\inv$ are continuous. Complete the following tasks:
	\begin{enumerate}[(a)]
		\ii Show that the interval $(0, 1)$ is homeomorphic to $\RR$.
		\ii Explain why, if $X$ and $Y$ are topologically equivalent
		metric spaces, then $X$ is \emph{connected} if and only if
		$Y$ is \emph{connected}.
		\ii Prove that the above statement is true
		if we substitute the word \emph{connected} with
		\emph{path connected} or \emph{compact}.
		\ii Show an example of a continuous bijective function
		$f\colon X\to Y$ that is not a homeomorphism.
		\ii Prove that if $X$ is compact then any
		continuous bijective function $f\colon X\to Y$ is a homeomorphism.
	\end{enumerate}
	
	We say that $X$ and $Y$ are \vocab{metrically equivalent}
	or \vocab{isometric} if there exists an \vocab{isometry}
	$f\colon X\to Y$, which is a bijective function that preserves distance,
	i.e. for all $x,\,y\in X$,
	\[d_Y(f(x), f(y)) = d_X(x, y).\]
	Show that
	\begin{enumerate}[(a)]
		\setcounter{enumi}{5}
		\ii if $f\colon X\to Y$ is an isometry then so is $f\inv$.
		\ii there exists a pair of metric spaces
		that are homeomorphic but not isometric.
	\end{enumerate}
\end{problem}

\begin{soln}
	First, for (a), it is easy to check that
	topological equivalence is an equivalence relation.
	Moreover, any two open intervals are homeomorphic because
	$g\colon (a, b)\to (c, d)$ given by
	\[g(x) = c + \frac{x-a}{b-a}\cdot (d-c)\]
	is easily checked to be a homeomorphism
	(it is continuous; it is bijective because it is monotonic and surjective;
	and its inverse is continuous because the inverse is also
	a linear polynomial in $x$).
	
	Now, consider the function $f\colon(-1,1)\to\RR$ given by 
	\[f(x) = \frac{x}{\sqrt{1 - x^2}}.\]
	This is clearly continuous, and we can check that
	it has an inverse
	\[f\inv(x) = \frac{x}{\sqrt{1 + x^2}}.\]
	The existence of the inverse implies that
	$f$ is bijective, and we can see that $f\inv$ is continuous.
	Thus $\RR$ is homeomorphic to $(-1,1)$, which is homeomorphic to $(0,1)$.\\
	
	For (b) and (c), we use the following fact
	(which justifies the phrase ``topologically equivalent''):
	\begin{theorem*}
		Let $f\colon X\to Y$ be a homeomorphism.
		Then $U\subseteq X$ is open if and only if $f(U)\subseteq Y$ is open.
	\end{theorem*}
	\begin{proof}
		First, if $f(U)$ is open in $Y$, then continuity says that
		$f\inv(f(U)) = U$ is open in $X$,
		where $f\inv(f(U)) = U$ because $f$ is bijective.
		
		Similarly, let $g = f\inv$; then, if $U$ is open in $X$,
		continuity of $g$ shows that $g\inv(U) = f(U)$ is open in $Y$.
	\end{proof}
	
	In each case, we only prove that
	$Y$ has the property if $X$ has it
	(the other direction follows by symmetry).\\
	
	For connectedness, we prove equivalently that
	$X$ is disconnected if and only if $Y$ is disconnected.
	Indeed, suppose that $X$ is a union of disjoint nonempty open sets
	$A$ and $B$. Then,
	\[f(A) = \{ f(x) : x\in A\},\qquad f(B) = \{f(x) : x\in B\}\]
	are open in $Y$ by the theorem and are nonempty.
	The union of these sets is
	\[f(A)\cup f(B) = \{f(x) : x\in A\text{ or }x\in B\}
	= \{f(x) : x\in A\cup B\} = \{f(x) : x\in X\} = Y\]
	with the last equality following from $f$ being a bijection.
	Similarly the intersection of the sets is
	\[f(A)\cap f(B) = \{f(x) : x\in A\text{ and }x\in B\}
	= \{f(x) : x\in A\cap B\} = \{f(x) : x\in\varnothing\}
	= \varnothing.\]
	Thus $f(A)$ and $f(B)$ are disjoint nonempty open sets
	whose union is $Y$, i.e. $Y$ is disconnected.\\
	
	For path connectedness, the idea is that
	we can push a path from one space to another
	because composition of continuous functions is continuous.
	Specifically, if $X$ is path connected, then let $a$ and $b$
	be points in $Y$. We can construct a path $\gamma\colon[0,1]\to X$
	between $f\inv(a)$ and $f\inv(b)$ so that $f\circ\gamma$
	is continuous and has $f(\gamma(0)) = a$ and $f(\gamma(1)) = b$,
	i.e. $f\circ\gamma$ is a path between $a$ and $b$.
	The other direction follows by the exact same logic.\\
	
	For compactness, let $\mathcal V = \{V_i : i\in I\}$
	be an open cover of $Y$. Then,
	\[\mathcal U = \{U_i : i\in I\}\]
	where $U_i = f\inv(V_i)$ is an open cover of $X$ because
	\[\bigcup_{i\in I} U_i =
	\bigcup_{i\in I} \{f\inv(y) : y\in V_i\} =
	\left\{f\inv(y) : y\in\bigcup_{i\in I} V_i\right\} =
	\{f\inv(y) : y\in Y\} = X\]
	with the last equality following from $f$ being bijective.
	Then, we can find a finite subcover of $X$, say
	$U_{i_1}$, \dots, $U_{i_n}$. Since $f$ is bijective,
	$f(U_{i_k}) = V_{i_k}$, and the union of this
	subcollection of $\mathcal V$ is
	\[\bigcup_{k=1}^n \{f(x) : x\in U_{i_k}\} =
	\left\{f(x) : x\in \bigcup_{k=1}^n U_{i_k}\right\} =
	\{f(x) : x\in X\} = Y,\]
	i.e. we have extracted a finite subcover of $Y$
	from $\mathcal V$.\\

	For (d), let $X = \NN$, $Y = \QQ$, and let $f\colon X\to Y$
	be any bijection between $\NN$ and $\QQ$.
	
	On one hand, $f$ is continuous because every point in $\NN$
	is an isolated point: to show that $f$ is continuous at
	$n\in\NN$, we can take $\delta < 1$ so that $B_X(n,\delta) = \{n\}$
	(because the only integer within $\delta$ of $n$ is $n$), and then
	\[|x - n| < \delta \iff x = n \implies |f(x) - f(n)| = 0 < \eps\]
	for any $\eps > 0$.
	
	On the other hand, $f\inv$ is not continuous at any point of $\QQ$.
	Taking $q\in\QQ$, we know because $f\inv$ is bijective that
	if $x\ne q$ is rational then $f(q)$ and $f(x)$
	are distinct elements of $\NN$, which differ by at least $1$.
	That is, regardless of the value of $|q - x|$,
	the inequality $|f(q) - f(x)| \ge 1$.
	Thus, taking $q_n = q + \frac 1n$, we have
	\[\lim_{n\to\infty} q_n = q,\]
	but the limit of $f(q_n)$ is not equal to $f(q)$ because
	$|f(q) - f(q_n)| \ge 1$ for all $n$,
	which finishes the proof of discontinuity.\\
	
	For (e), to show that $f\inv$ is continuous, we can show that
	$f$ maps closed sets to closed sets. This is pretty simple:
	if $E$ is a closed subset of $X$, then it is also compact.
	Then, $f(E)$ is also compact, hence closed!\\
	
	For (f), we have that, given $x$ and $y$ in $Y$,
	we can write $x = f(u)$, $y = f(v)$ where $u$, $v$ are in $X$. Then
	\[d_Y(x, y) = d_Y(f(u), f(v)) = d_X(u, v) = d_X(f\inv(x), f\inv(y)).\]
	
	For (g), observe that $(0, 1)$ and $\RR$ are homeomorphic
	(as shown in (a)), but they are not isometric.
	Indeed, if they were isometric, then there exist $\alpha$ and $\beta$
	in $(0, 1)$ such that $f(\alpha) = 0$ and $f(\beta) = 2$, so that
	\[2 = |f(\alpha) - f(\beta)| = |\alpha - \beta|\]
	which is a contradiction because $\operatorname{diam}((0,1)) = 1 < 2$.
\end{soln}

\begin{problem}
	Let $\mb a = (a_1, a_2, \dots)$ be an element of $\ell^2$.
	Show that the function $\varphi_{\mb a}\colon\ell^2\to\RR$ defined by
	\[\varphi_{\mb a}(\mb x) = \sum_{i=1}^\infty a_ix_i\]
	is well-defined (i.e. the series on the right converges)
	and is linear and continuous.
	
	In 1908, the Hungarian mathematician Frigyes Riesz proved that
	any continuous linear function $\ell^2\to\RR$
	is expressible in this way --- try proving this!
\end{problem}

\begin{soln}
	The series converges because it can be bounded termwise by
	\[|a_ix_i| \le \frac{a_i^2 + x_i^2}{2}\]
	(by expanding and reducing to the equivalent $(|a_i| - |x_i|)^2 \ge 0$).
	This shows that $\sum a_ix_i$ converges absolutely, so it converges.
	
	It is linear because
	\[\varphi_{\mb a}(\mb x) + \varphi_{\mb a}(\mb y)
	= \sum_{i=1}^\infty a_ix_i + \sum_{i=1}^\infty a_iy_i
	= \sum_{i=1}^\infty a_i(x_i + y_i)
	= \varphi_{\mb a}(\mb x + \mb y).\]
		
	To finish, we have
	\[\varphi_{\mb a}(\mb x) \le \frac{\|\mb a\|^2 + \|\mb x\|^2}{2},\]
	so, when $\mb x$ has norm at most $1$,
	\[\varphi_{\mb a}(\mb x) \le \frac{\|\mb a\|^2 + 1}{2},\]
	hence $\varphi_{\mb a}$ is bounded on the closed unit ball;
	Problem 13 then allows us to conclude that $\varphi_{\mb a}$
	is continuous.
\end{soln}














\end{document}